\newpage
\section{Introduction and Cloud Paradigm}

\subsection{Institutional Context and Operational Demands}
Higher educational institutions such as the Indian Institute of Information Technology Design and Manufacturing Kurnool (IIITDM Kurnool) operate in complex, data-centric environments requiring synchronized management of curriculum, student admissions, course enrollments, attendance tracking, assignment evaluation, and official academic grade reporting \cite{iiitdm_kurnool}. Academic operations demand reliable, low-latency access across thousands of concurrent student, faculty, and administrative users during peak evaluation cycles.

\subsection{Paradigm Shift from Legacy Infrastructure to Cloud Computing}
Historically, educational management software operated on physical, on-premise servers located within campus data centers. While functional, on-premise deployments introduce severe architectural constraints:
\begin{itemize}[leftmargin=*,itemsep=2pt]
    \item \textbf{Inflexible Capacity and High CapEx}: On-premise hardware must be provisioned for peak annual load (e.g., course registration or result publication), leaving servers underutilized during standard instructional periods.
    \item \textbf{Single Points of Failure (SPOF)}: Local server outages, power interruptions, or network partition events lead to complete operational blackout without automated cloud failover.
    \item \textbf{Administrative Burden}: Operating system patches, manual database backups, disk storage expansion, and hardware lifecycle refreshes consume significant institutional engineering effort.
\end{itemize}

\subsection{AWS Cloud Computing Foundation}
To overcome these limitations, the \textbf{CloudCampus} system is architected natively on \textbf{Amazon Web Services (AWS)} in Region \texttt{us-east-1} \cite{aws_ec2, aws_s3}. Cloud computing provides an elastic, utility-based delivery model where compute instances, managed database engines, object storage buckets, and serverless functions are provisioned on-demand, strictly isolated through Virtual Private Clouds (VPCs), and governed through identity federation and programmatic security policies.
